%%%%%%%%%%%%%%%%%%%%%%%%%%%%%
%%                          %%
%% Formát a jazyk dokumentu %%
%%                          %%
%%%%%%%%%%%%%%%%%%%%%%%%%%%%%%
\chapter{Formát a jazyk}\label{sec:formatLanguage}
% TODO: Treba upravit nizsie
\section{Formát dokumentu}
Rozmery stránky, typy písma, veľkosti, riadkovanie,
medzery medzi odsekmi, formát nadpisov, obrázkov, tabuliek,
rovníc a~ďalšie vizuálne parametre záverečnej práce
rešpektujú do maximálnej miery normu STN 01 6910: 2023
Pravidlá písania a úpravy písomností~\cite{stn016910}.

\subsection*{Rozmery strany}
Veľkosť bežnej textovej strany záverečnej práce je A4,
t.~j. 21\,cm $\times$ 29,7\,cm.
Pravý a~ľavý okraj majú šírku 2,75\,cm,
horný a~dolný okraj majú výšku 3\,cm.
Päta stránky, v~ktorej sa nachádza číslo strany,
je od spodnej hrany stránky vzdialená o~1,25\,cm.
Šírka textu je 15,5\,cm, jeho výška 23,7\,cm.
Horný a~dolný okraj obálky sú z~estetických
dôvodov zmenšené na 2\,cm.

\subsection*{Písmo a riadkovanie}
Základný font šablóny je normálny rez tzv. antikvového písma
s~veľkosťou 12\,pt.
V~tejto šablóne je to Computer Modern.
Vhodné sú aj iné fonty s pätkami ako Times, Georgia, Palatino a~podobne.
Na obálke a~titulnom liste používame bezpätkový (grotesk) font Latin Modern.
Jednotlivé typy odsekov (nadpisy, poznámky a~pod.)
majú jednotný typ písma,
odlišnosti vyjadrujeme rezom (polotučné písmo, kurzíva)
alebo veľkosťou.

Parameter \verb|\linespread| má hodnotu 1,25, t.~j.
vzdialenosť riadkov textu vo veľkosti 12\,pt je 15,6\,pt.

\subsection*{Nadpisy}
Šablóna záverečnej práce FEIstyle je založená na
štandardnej šablóne \LaTeX-u article.
Nadpis najvyššej úrovne je \verb|\section| zodpovedajúci kapitole.
Podkapitoly sú \verb|\subsection| a~\verb|\subsubsection|.
Číslovanie kapitol a podkapitol je viacúrovňové typu X.Y.Z, kde X je číslo kapitoly, Y je číslo podkapitoly a Z je číslo časti podkapitoly.
Číslovanie vyšších úrovní nie je definované.
Tvar a forma nadpisov zodpovedá norme STN ISO 2145: 1978 Dokumentácia.
Číslovanie oddielov a~pododdielov písaných dokumentov~\cite{iso2145}.

Nová kapitola začína vždy na novej strane.
Príkaz \verb|\section| spôsobí okrem sadzby čísla a názvu kapitoly aj ukončenie predošlej kapitoly, vysádzanie všetkých plávajúcich objektov (obrázky, tabuľky, výpisy kódu), ktoré sa nepodarilo umiestniť na príslušné miesto v texte, a prejde na novú stranu.

\section{Jazyk a gramatika}
Záverečná práca na FEI STU v~Bratislave musí byť napísaná
buď po slovensky alebo po anglicky.
Ak je jazyk práce angličtina, musí po závere nasledovať
rezumé v~slovenskom jazyku.

Záverečná práca univerzitného štúdia sa vyznačuje
vysokou jazykovou úrovňou.
Gramatické a štylistické chyby sú neprípustné.
Študent by mal tejto stránke diela venovať patričnú
pozornosť a~podľa možností nechať rukopis prejsť
kvalifikovanou jazykovou kontrolou.
Najmä bakalárska práca predstavuje v~živote väčšiny študentov
prvý rozsiahlejší autorský útvar,
ktorý má významný vplyv na jeho ďalší život a~kariéru.

Aj keď väčšina textových editorov dokáže odhaľovať preklepy,
neporadí si s~komplikovanejšou gramatikou a~štylistikou.
Treba sa riadiť najmä pravidlami slovenského pravopisu,
slovníkmi slovenského jazyka a~ďalšími zdrojmi,
ktoré možno nájsť na webových stránkach
Jazykovedného ústavu Ľudovíta Štúra SAV.%
\footnote{\href{https://www.juls.savba.sk/}{\tt www.juls.savba.sk}}
Využiť môžeme aj jazykovú poradňu,
ktorú poskytuje ústav bezplatne a~to buď telefonicky alebo
prostredníctvom emailovej komunikácie.
Cenným zdrojom informácií môže byť aj Jazyková poradňa
denníka SME v spolupráci
s~Jazykovedným ústavom Ľudovíta Štúra SAV%
\footnote{\href{https://jazykovaporadna.sme.sk/}{\tt jazykovaporadna.sme.sk}}
alebo online slovníky slovenského jazyka,%
\footnote{\href{https://slovnik.juls.savba.sk/}{\tt slovnik.juls.savba.sk}}
prípadne národný jazykový korpus.%
\footnote{\href{https://korpus.sk/}{\tt korpus.sk}}

Pri písaní práce dbáme najmä na pravopisné javy ako sú písanie
tvrdého a~mäkkého y/i vo vybraných slovách,
v~príponách a~koncovkách pri skloňovaní
(pekný muž, ale pekní muži),
v~číslovkách (rozprávali sme sa so siedmimi v~poradí
-- skončili siedmi v~poradí,
ale hrali sme sa so siedmymi deťmi -- detí bolo sedem), atď.
Rovnako dôležité je správne písanie rodov,
skloňovanie a časovanie.

Veľmi komplexná a~dôležitá zložka gramatiky
je písanie čiarok v~súvetiach.

Popri gramatike je podstatná aj štylistická tvorba viet,
ktorú musí študent univerzity zvládať na vysokej úrovni.

\subsection{Delenie slov}
Tzv. \emph{textové procesory} ako MS Word, LibreOffice a~Apache OpenOffice
ponúkajú automatické delenie slov na konci riadka.
Systém na sadzbu textu \LaTeX\ má túto funkciu automaticky zapnutú a~jej slovenská lokalizácia je veľmi kvalitne spracovaná.

Vo veľkej väčšine prípadov je delenie
v~súlade s~pravidlami jazyka.
Môžu sa vyskytnúť sporné okolnosti,
kedy počítač nerozdelí slovo správne.
Väčšinou máme možnosť do procesu zasiahnuť
a~ručne kontrolovať delenie slov na miestach,
s ktorými si softvér nevie poradiť.
Príkaz na preferované rozdelenie slova je \verb|\-|.
Napríklad slovo \verb|predstave\-nie| \LaTeX\
preferovane rozdelí v mieste prípony.

V každom prípade je žiadúce slová na konci riadka deliť
a\ túto možnosť nevypínať.
Prospieva to práci ako po technickej,
tak aj po estetickej stránke.
Odseky obsahujú menej dier,
textová oblasť stránky je vyplnená homogénnejšie,
čo prispieva k~lepšej čitateľnosti.
V prípade, že používame zarovnávanie do bloku tak,
ako aj v~tomto dokumente,
je prítomnosť dier v~odseku značne rušivá.
Ak používame zarovnanie textu doľava,
nepoužívanie delenia slov má vplyv na vznik tzv. riek,
čo je náhle striedanie dlhých a~krátkych riadkov.
Pravý okraj textu je nepekne zubatý.

Pravidlá rozdeľovania slov na konci riadka sú pomerne zložité.
Základné pravidlo, ktoré si pamätáme zo základnej školy,
je, že slová delíme na slabiky pred spoluhláskou alebo medzi
dvomi spoluhláskami.
Ak si nie sme istí, uprednostňujeme delenie v mieste,
kde sa ku koreňu slova pripájajú predpony alebo prípony,
prípadne v~mieste spojenia slov v~zloženom slove.

Pri slovách utvorených predponou alebo príponou
uprednostňujeme morfologické delenie
pred rozdelením koreňa slova.
Najskôr sa snažíme deliť slovo za predponou,
ak to nejde, skúsime to pred príponou.
Napríklad slovo \emph{predstavenie}
delíme na slabiky takto: \emph{pred-sta-ve-nie}.
Pri rozdeľovaní slov uprednostňujeme model
\emph{pred-stavenie}, výnimočne aj \emph{pred-stave-nie}.
V slove \emph{výklenok} sa uplatňuje pravidlo morfologického
delenia pred delením v~mieste zhluku spoluhlások.
Sylabická stavba tohto slova je \emph{vý-kle-nok},
nie \emph{výk-le-nok},
pretože slovo pozostáva z troch častí: predpony \emph{vý},
koreňa \emph{kle} a prípony \emph{nok}.
Mohli by sme namietať, že prípona je \emph{ok},
pomocou ktorej bolo vytvorené podstatné meno zo slovesa
klenúť alebo z~prídavného mena klenutý,
kde identifikujeme koreň \emph{klen}.
V skutočnosti je však príponou \emph{-nok}.
Morfológia je pomerne komplexná problematika,
a~nedokážeme tu obsiahnuť všetky jej detaily.
Väčšinou sa môžeme spoľahnúť na softvér,
že slová rozdelí správne.
V prípade pochybností využijeme externé pomôcky spomenuté
v~úvode tejto kapitoly.

Slová spojené spojovníkom rozdeľujeme v~mieste spojovníka tak,
že spojovník napíšeme na konci aj na začiatku riadka.
Slovo \emph{vedecko-pedagogický} môžeme rozdeliť takto:
\emph{ve-dec-ko-pe-da-go-gic-ký}.
Ak delenie padne na miesto spojenia slov,
rozdelíme ho nasledujúcim spôsobom:\\
\indent\emph{vedecko-}\\
\indent\emph{-pedagogický}

V šablóne rieši tento problém príkaz
\verb|\languageattribute{slovak}{split}|,
ktorý je súčasťou jazykového balíka \verb|babel|.

Nesprávne delenie slov sa v~práci zvyčajne objaví
len zriedkavo a~nemá vplyv na jej hodnotenie.
Netreba sa naň príliš sústrediť a~robiť si starosti.
Celkový vzhľad práce viac naruší vypnutie delenia slov,
než občasná malá chyba.

\subsection{Jednopísmenové predložky a spojky}
Hovoríme o predložkách k, o, v, s, z,
ktoré by nemali ostať osamotené na konci riadka.
Do tejto kategórie patria aj spojky a, i.
Jednopísmenové slová pripájame k~nasledujúcemu slovu pomocou
tzv. \emph{nedeliteľnej medzery},
čo je špeciálny netlačiteľný znak.
V kódovaní UTF-8 má číslo 00A0 (ASCII 160)
a~hovorí textovému procesoru,
že na tomto mieste nesmie byť za žiadnych okolností
koniec riadka.
V programe MS Word ho zadáme použitím klávesovej skratky
Ctrl-Shift-Medzera.
V \LaTeX-u\ zapíšeme nedeliteľnú medzeru s~premenlivou šírkou
ako symbol vlnovka (\verb|~|)
Napríklad slovné spojenie \emph{v~priestore} napíšeme takto:
\verb|v~priestore|.
Na rozdiel od Wordu, \LaTeX\ takéto medzery nevkladá automaticky
a treba to urobiť ručne.

Existuje viacero medzier, ktoré sú tiež nedeliteľné a~majú pevnú šírku.
Najpoužívanejšia tzv. úzka medzera a zapíšeme ju ako \verb|\,|.
Takýto typ medzery používame pri zápise hodnôt fyzikálnych veličín a vkladáme ju medzi číslo a jednotku.

\section{Štylistika}
Niektorí oponenti vyčítajú študentom príliš dlhé súvetia,
iní zas príliš krátke.
Pravda je, že jednoduché vety pôsobia školácky,
zatiaľ čo dlhé súvetia sú často nezrozumiteľné a~únavné.

V prvom rade sa snažíme nevrstviť podraďovacie súvetia.
Vo vete \emph{Elektrostatické pole je fyzikálne pole,
ktoré tvoria elektrické náboje, ktoré sú v pokoji}
je dvakrát použitá spojka ktoré,
čo je síce prípustné, avšak nie príliš estetické.
Vetu môžeme opraviť takto:
\emph{Elektrostatické pole je fyzikálne
pole tvorené elektrickými nábojmi, ktoré sú v pokoji.}
Ak sa chceme vyhnúť trpnému rodu,
môžeme vetu preformulovať nasledujúcim spôsobom:
\emph{Elektrostatické pole tvoria elektrické náboje,
ktoré sú v~pokoji.}
Vypadol síce pojem fyzikálne pole,
ale zmysel vety zostal nezmenený.

Správne a~plynulo bude veta vyzerať aj v~tomto tvare:
\emph{Elektrostatické pole je fyzikálne pole,
ktoré tvoria elektrické náboje v pokoji.}
V prípade potreby môžeme vetu napísať aj inak:
\emph{Fyzikálne pole elektrických nábojov,
ktoré sú v pokoji, nazývame elektrostatické pole.}

Obmieňame štruktúru po sebe nasledujúcich viet:
\emph{Z výsledkov merania je zrejmé,
že predpoklad o~zvyšovaní pohyblivosti nosičov náboja
s~teplotou bol správny.
Na začiatku práce sme hovorili o~tom,
že toto tvrdenie podporíme hodnovernými experimentálnymi dátami.}
Obe súvetia sú podraďovacie so spojkou že.
Aby sme sa vyhli opakovaniu rovnakého typu viet,
môžeme prvú vetu prepísať:
\emph{Výsledky merania potvrdili predpoklad o~zvyšovaní
pohyblivosti nosičov náboja s~rastúcou teplotou.}
Druhú vetu ponecháme bez zmeny.

Veľmi osviežujúco pôsobí, ak medzi dlhé a~kvetnaté súvetia
občas vložíme jednoduchú holú vetu.
Použijeme predchádzajúci príklad:
\emph{Na začiatku práce sme hovorili o~tom,
že predpoklad o~zvyšovaní pohyblivosti nosičov náboja
s~rastúcou teplotou podporíme hodnovernými
experimentálnymi dátami.
Merania ho potvrdili.}
Tento malý trik je nečakane účinný a~prispieva k~lepšiemu
toku myšlienok.

Pozor, v~texte pozostávajúcom z~krátkych jednoduchých viet
je niekoľkoriadkové súvetie desivé:
\emph{Pohyblivosť rastie s~teplotou.
Hovorili sme o~tom už na začiatku.
Tvrdenie ešte podporíme experimentom.
Ukazuje sa, že sme predpoklad o~rastúcej pohyblivosti
nosičov náboja so zvyšujúcou sa teplotou,
pokiaľ berieme do úvahy výsledky meraní,
formulovali správne.}

Aby bol písaný text zaujímavý a~udržal čitateľov záujem,
používame stredne dlhé súvetia pozostávajúce maximálne
z~dvoch až troch viet.
Občas text oživíme jednoduchou krátkou vetou.
Dávame si pri tom pozor,
aby táto činnosť nebola príliš schematická.

\section{Anglický jazyk}
Šablóna FEIstyle podporuje slovenský a anglický jazyk.
Pre prácu v anglickom jazyku je potrebné túto skutočnosť nastaviť v preambule hlavného súboru \verb|thesis.tex| ako nepovinný parameter \verb|en| príkazu definície šablóny
\begin{trivlist}
\item \texttt{\textbackslash documentclass[bp,en]\{FEIstyle\}}
\end{trivlist}
Prvý parameter určuje, či ide o bakalársku (\verb|bp|) alebo diplomovú (\verb|dp|) prácu.

Anglická práca musí obsahovať po závere rezumé a na to je potrebné odstrániť komentár pred príkazom
\verb|\FEIresume{includes/resume}|.

Na jazykovú lokalizáciu používame balík \verb|babel|.
Ak sa v práci písanej v slovenčine nachádzajú výrazy v angličtine,
uvedieme takýto text do makra \verb|\foreignlanguage|, čím zabezpečíme správne medzery a~delenie slov.
Napríklad pri zavádzaní skratky AI môžeme napísať, že ide o anglický výraz pre umelú inteligencie \foreignlanguage{english}{Artificial Intelligence}.
Zapíšeme ho nasledujúcim spôsobom:
\begin{verbatim}
\foreignlanguage{english}{Artificial Intelligence}
\end{verbatim}

Ak nastavíme parametrom \verb|en| anglický jazyk ako hlavný, stane sa slovenčina cudzím jazykom.

\section{Použitie umelej inteligencie}\label{sec:utilizingAI}
Na optimalizáciu formulácie myšlienok môžeme využiť služby
umelej inteligencie (AI, z~ang. \emph{\foreignlanguage{english}{artificial intelligence}}) a~tzv. veľkých jazykových modelov (LLM, z~ang. \emph{\foreignlanguage{english}{large language model}}).
Umelá inteligencia dokáže kontrolovať rozsiahlejšie časti prác,
vyhľadáva chyby a navrhuje vhodnejšie
formulácie na základe pravidiel,
ktoré sme aplikovali v predchádzajúcom texte.
Neosvedčuje sa však pri kompozícii textov.
Neuspokojivé výsledky dosahujeme aj v prípadoch,
kedy necháme umelú inteligenciu preformulovať celé odseky.
Zanáša do nich chyby a nezmysly, ktoré tam pôvodne neboli.
Ťažko sa potom odhaľujú.
Tento jav poznáme ako tzv.
halucinácie a trpia nimi všetky nástroje AI,
vrátane najznámejšieho ChatGPT.

Napriek tomu predstavujú služby AI silný nástroj pri tvorbe pôvodného obsahu, zvlášť užitočné sú tzv. generatívne umelé inteligencie (GAI),
ktoré dokážu vytvárať výstupy takmer na nerozoznanie od tvorby človeka.
Ich správna aplikácia nepochybne prispieva k~vyššej jazykovej a obsahovej kvalite záverečných prác.
Treba však mať na pamäti, že záverečná práca má byť
pôvodné autorské dielo študenta a všetky časti,
ktoré nepochádzajú od autora musia byť riadne
zdokumentované a deklarované v~zozname použitých zdrojov.
V žiadnom prípade sa neodporúča, aby
GAI formulovala pôvodné myšlienky
alebo súvislé časti práce.
Takéto konanie považujeme za nečestné podobne,
ako keby prácu písal niekto iný,
prípadne by boli celé odseky prebrané z iného zdroja bez korektného citovania (pozri kapitolu \ref{sec:citation}).

Používane umelej inteligencie pri písaní záverečných prác
upravuje opatrenie rektora STU v Bratislave č. 1/2024-O, ktoré budeme ďalej v~texte uvádzať ako \uv{opatrenie}~\cite{opatrenie12024}.

\subsubsection*{Povolené činnosti umelej inteligencie bez potreby deklarácie}
Podľa čl. V, ods. 2, písm. a) opatrenia môžu študenti používať GAI bez potreby deklarácie na tieto činnosti: kontrola gramatiky, oprava textu, tvorba osnovy, zhromažďovanie informácií a~použitie výpočtových metód a~softvérov, ktoré obsahujú prvky AI.

\subsubsection*{Deklarácia činnosti generatívnej umelej inteligencie}
Čl. V, ods. 2, písmeno b) opatrenia obsahuje zoznam možností použitia GAI, ktoré je potrebné v práci deklarovať na konci po zozname literatúry.
Ide o nasledujúce činnosti: preklady medzi jazykmi, úpravy a reformulácie textu, tvorba zhrnutia a rešerší, citovanie odpovedí GAI, tvorba počítačových programov, tvorba grafického obsahu a obrázkov.

V závere práce, uvedieme za zoznamom literatúry
časti textu vytvorené s~pomocou AI, spôsob ich využitia a použitý nástroj AI (Opatrenie \cite{opatrenie12024}, čl. VI., ods. 2).

V hlavnom súbore záverečnej práce \verb|thesis.tex| je príkaz
\verb|\FEIaiDeclaration| na načítanie súboru \verb|ai_declaration.tex| z~priečinka \verb|includes|.
Každý výskyt použitia nástrojov AI zapíšeme ako
položku \verb|\item| do pripraveného prostredia
\verb|trivlist| v tomto súbore.
Formát a~obsah jednotlivých záznamov je naznačený v prílohe opatrenia číslo 1/2024-O.
Záznamy obsahujú tieto prvky:
\begin{trivlist}
  \item Názov spoločnosti (dátum), Názov nástroja, čast práce, účel použitia.
\end{trivlist}

Predchádzajúci vzorec vygeneroval nástroj ChatGPT 4o od firmy
OpenAI dňa 2. 2. 2025 na základe analýzy spomínaného opatrenia.
V deklarácii použitia umelej inteligencie sa zapíšeme tento
záznam:
\begin{trivlist}
  \item OpenAI (2025), ChatGPT 4o, časť \ref{sec:utilizingAI},
  generovanie vzorca záznamu použitia AI.
\end{trivlist}

Súčasná verzia šablóny FEIstyle nedisponuje nástrojmi na
automatizáciu záznamov činnosti AI.
Preto ich treba zapisovať ručne do súboru
\verb|includes/ai_declaration.tex|.
